\section{Models, oracles, algorithms, and theorems}\label{sec:statements}

The objective throughout is the literal inequality
\begin{equation}\label{eq:target}
 \E\ell(\widehat z,z)\le C R_*(n),\qquad C<\infty\text{ independent of }n.
\end{equation}
An estimate of the form $\exp\{-(1-o(1))I_n\}$ does not establish
\eqref{eq:target} when the oracle error is of order
$d^{-1/2}\exp(-I_n)$. The proof below controls bounded changes in a
tilted cumulant and retains that prefactor. For comparison, the main theorem
of Zhang~\cite{zhang} characterizes the leading error exponent of a spectral
clustering algorithm, including a polynomial exceptional-event term. We use
its distinction between an actual spectral algorithm and its population
surrogate, while proving a different likelihood comparison on the parameter
ranges stated here.

\subsection{Bernoulli models and fixed structural constants}

There are $n$ vertices and $K\ge2$ groups, with fixed $K$.
Let $z\in[K]^n$, let $Z$ be its indicator matrix, and let $n_a$ be the size
of group $a$. Throughout $n_a\ge\pi_-n$ for a fixed $\pi_->0$.
The edges above the diagonal are independent, the adjacency matrix $A$ is
symmetric, and $A_{ii}=0$. In DC-SBM,
\begin{equation}\label{eq:mean-model}
 A_{ij}\sim\Ber(\theta_i\theta_j p C_{z_i z_j}),\quad i<j,
 \qquad d=np,\qquad
 \Omega=\Theta Z(pC)Z^\top\Theta.
\end{equation}
Here $\Theta=\diag(\theta_i)$, $C=C^\top$, and the fixed structural bounds are
\begin{equation}\label{eq:structural}
 0<\theta_-\le\theta_i\le\theta_+<\infty,\quad
 \sum_{i:z_i=a}\theta_i=n_a,\quad
 0<c_-\le C_{ab}\le c_+<\infty,\quad
 \sigma_{\min}(C)\ge\kappa>0.
\end{equation}
The model class is nonempty only when $1\in[\theta_-,\theta_+]$.
All constants in asymptotic bounds may depend on these structural bounds,
but not on $n$. One may take a fixed $C$, or a sequence in this same compact
class. Ordinary SBM is the special case $\theta_i=1$, with block probability
matrix $P=pC$.

The completed matrix $\Omega$ includes a diagonal for algebraic convenience;
the actual mean is $\E A=\Omega-\diag(\Omega_{ii})$.
That deterministic difference has norm $O(p)$ and is included in all error
bounds. Put $D=\diag(A\one)$ and
$\lambda_i=(\E A\one)_i$. Thus $\lambda_i\asymp d$ and
$\Omega\one$ equals $d\theta_i S_{z_i}$ in coordinate $i$, where
\begin{equation}\label{eq:exposures}
 \pi_a=n_a/n,\qquad S_a=\sum_b\pi_b C_{ab}.
\end{equation}
The loss for the estimated partition is
\[
 \ell(\widehat z,z)=\frac1n\min_{\sigma\in\mathfrak S_K}
 \sum_i\ind{\sigma(\widehat z_i)\ne z_i}.
\]

\subsection{Three explicitly different local oracles}\label{sec:oracle}

All oracle classifiers below use the raw incident row $A_{i,-i}$ and the
true neighbor labels. They are local likelihood classifiers: an exact global
block-size constraint is not used to determine the one missing label.
For simplicity candidate weights are uniform. Any fixed weights bounded
above and below by positive constants give the same conclusions.

For ordinary SBM, the oracle chooses the maximizer of
\begin{equation}\label{eq:sbm-oracle}
 \ell_{i,*}(a)=\sum_{j\ne i}
 \{A_{ij}\log P_{a z_j}+(1-A_{ij})\log(1-P_{a z_j})\}.
\end{equation}
Its average error probability is $R_{\mathrm{SBM},*}$.
For DC-SBM, the \emph{known-target-degree oracle} replaces $P_{a z_j}$ in
\eqref{eq:sbm-oracle} by $\theta_i\theta_j p C_{a z_j}$ and is supplied
with every true $\theta_j$, including $\theta_i$. Its risk is
$R_{\theta,*}$.

The \emph{profile oracle} is supplied with true neighbor parameters and
$pC$, but profiles the target degree separately for every candidate:
\begin{equation}\label{eq:profile-oracle}
 \ell_{i,\mathrm{prof},*}(a)=
 \sup_{t\ge0}\sum_{j\ne i}
 \{A_{ij}\log(t\theta_jpC_{a z_j})
 +(1-A_{ij})\log(1-t\theta_jpC_{a z_j})\}.
\end{equation}
We restrict candidate edge probabilities to at most $q_{\max}=1/2$; the
constraint is inactive on the events used in the sparse theorem. Write
$R_{\mathrm{prof},*}$ for its average error probability. This is a
full-edge generalized likelihood-ratio reference. It is not described as a
Bayes test between two fixed known-$\theta_i$ laws. In all three definitions,
the risk averages the actual directional error at each fixed true vertex.
Ties are resolved by a fixed ordering; strict-tail arguments below cover
any such ordering.

\subsection{Primary result: no logarithmic degree lower bound}

\begin{theorem}[Regularized fitted Bernoulli likelihood at every diverging degree scale]
\label{thm:regularized}
Assume \eqref{eq:mean-model}--\eqref{eq:structural}, with fixed arbitrary
$K$, and suppose
\[
 d=np\longrightarrow\infty,\qquad np^2\le C_1<\infty.
\]
Use the empirical degree-floor encoder and finite block-LRT decoder
specified in \cref{restore:algorithm}. It retains only the raw degrees
and the $K$ informative eigenpairs of $G^{-1/2}AG^{-1/2}$, where
$G_i=\max(D_i,a_0\bar D)$ and $a_0>0$ is fixed and sufficiently small.
After exact-degree reconstruction and hollowing, use radius-count
initialization and at least one block Bernoulli fit-and-LRT update.
Any finite sequence of the specified hard refits, including finite
adaptive stopping, satisfies
\[
 \E\ell(\widehat z,z)\le
 \begin{cases}
 C R_{\mathrm{SBM},*}(n),&
    \theta_i=1\text{ known, with no target profiling},\\
 C R_{\mathrm{prof},*}(n),&
    \text{DC-SBM with the specified nuisance fit and target profiling}.
 \end{cases}
\]
Here $C$ is independent of $n$ on the fixed structural class.
Under the finite-$n$ equal-exposure condition
\eqref{eq:equal-exposure}, the DC bound is also
$C R_{\theta,*}(n)$. No lower bound proportional to $\log n$,
correct initial labels, true nuisance parameters, sample splitting,
or deleted-graph eigensystem is an input to this algorithm.
\end{theorem}

The complete algorithm and proof are in \cref{sec:sparse-restored},
especially \cref{restore:dc-theorem,restore:sbm-theorem}.
This result permits $np$ to diverge arbitrarily slowly. The
logarithmic scale appears only in the choice of probability estimate
within its proof. Ordinary SBM uses a block Bernoulli fit here;
the distinct free-profile mixture schedule below retains its own
separately proved scope.

\subsection{Earlier unregularized encoders and their separate results}

The rest of this section records the unregularized algorithms analyzed
in the earlier sections of the note. Their logarithmic-range hypotheses
are not hypotheses of \cref{thm:regularized}. Their positive-degree
fallback convention is replaced by the explicit floor in the primary
algorithm.

When all degrees are positive, form
\[
 S=D^{-1/2}AD^{-1/2}=I-L.
\]
Retain the degree eigenvector $D^{1/2}\one/\norm{D^{1/2}\one}_2$ and the
$K-1$ other eigenvectors with largest absolute eigenvalues in its orthogonal
complement. Let $U^\top U=I_K$ and $SU=U\Lambda$, where
$\Lambda=\diag(\mu_1,\ldots,\mu_K)$. On the proved event this is exactly
the informative $K$-dimensional space. The encoder supplies only
$(D,U,\Lambda)$ to the decoder. Put
\begin{equation}\label{eq:encoder-action}
 X=D^{-1/2}U,\qquad AX=D^{1/2}U\Lambda.
\end{equation}
Thus row actions on the retained space require no later edge lookup.
The normalized Laplacian eigenvalues contain the same information via
$\mu=1-\lambda_L$. Exact arithmetic and algebraic retained spaces define
the theorem; a fixed numerical rank threshold is not its definition.

The ordinary-SBM decoder uses the \emph{signed} reconstruction
\begin{equation}\label{eq:euclidean-decoder}
 Q=\operatorname{orth}(X),\quad \Pi_X=QQ^\top,\quad R_E=A\Pi_X.
\end{equation}
It fits free Bernoulli row profiles with the finite growing/replacement
schedule and mandatory final E/M step specified in
\cref{sec:sbm-proof}. Only probabilities are clipped; $R_E$ is not clipped.
Neither true labels nor a correct initial partition are inputs.

For DC-SBM we make the explicit degree-metric correction
\begin{equation}\label{eq:metric-decoder}
 H_D=D^{-1/2}UU^\top D^{1/2},\qquad
 R_D=AH_D=D^{1/2}U\Lambda U^\top D^{1/2}.
\end{equation}
With $h_i=\norm{U_{i\cdot}}_2^2$ and
$s_i=(R_D)_{ii}/(1-h_i)$, its hollow version is
\begin{equation}\label{eq:hollow-decoder}
 R^\circ=R_D+\diag(s)(H_D-I).
\end{equation}
It satisfies
$\diag(R^\circ)=0$, $R^\circ\one=D\one$ and $R^\circ X=AX$.
It need not be symmetric after hollowing; a symmetric block parameter is
fitted with the ordered-row working criterion divided by two.

For the initial partition, normalize every row of $U$ to unit length.
Repeatedly select a remaining row having the largest number of remaining
rows within distance $1/8$, and remove rows within distance $1/4$ of the
selected row. After $K$ selections assign every row to the closest selected
row. The separation argument in \cref{sec:dcsbm-proof} proves its
$O(1/d)$ error; this is a deterministic procedure with $O(Kn^2)$ distance
comparisons, not an assumed global k-means solution.
Normalize degree estimates to mean one within each fitted group, fit
the symmetric block matrix by the exact Bernoulli working criterion, and
perform at least one exact Bernoulli profile-LRT assignment using
$R^\circ$. One pass suffices. The same theorem permits a finite number of
the stated hard refits, including adaptive stopping, by the invariant-basin
argument proved below.

Both algorithms return a fixed partition, or their initial partition,
if a prescribed operation is undefined. This includes zero degrees,
insufficient retained rank, a hollow denominator zero, an empty initial
group, or an inadmissible parameter fit. The required operations are all
well-defined on the good event, so this convention costs only its actual
failure probability.

\begin{theorem}[Earlier ordinary free-profile decoder]\label{thm:sbm}
Fix $K\ge2$, $\pi_->0$ and $\kappa>0$. Suppose
\[
 n_a\ge\pi_-n,\qquad
 p/2\le P_{ab}\le p\le1/4,\qquad
 \sigma_{\min}(P)\ge\kappa p,\qquad d=np\ge5\log n.
\]
Use \eqref{eq:euclidean-decoder} and the finite growing/replacement
likelihood decoder, with probability floor
$\varepsilon=c_0\max(1,n^{-1}\sum_iD_i)/n$ for a fixed sufficiently small
$c_0>0$ (for example any $c_0<1/16$ on this class).
There are constants $C,n_0$, depending only on the displayed fixed
structural constants and $c_0$, such that for $n\ge n_0$,
\begin{equation}\label{eq:sbm-risk}
 \E\ell(\widehat z_E,z)\le C R_{\mathrm{SBM},*}(n).
\end{equation}
The algorithm uses neither $P$, $p$, nor the true partition.
This includes fixed arbitrary $K$, unbalanced group sizes of comparable
order, and general symmetric full-rank block probability matrices in the
displayed range. It does not require all within-group probabilities, or all
between-group probabilities, to be identical.
\end{theorem}

\begin{theorem}[Earlier unregularized DC profile decoder]
\label{thm:dc}
Assume \eqref{eq:mean-model}--\eqref{eq:structural}. Fix $C_1<\infty$ and
suppose
\begin{equation}\label{eq:dc-range}
 C_0\log n\le d=np,\qquad np^2\le C_1,
\end{equation}
where $C_0$ is sufficiently large depending only on the fixed structural
bounds and $C_1$. The degree-metric decoder
\eqref{eq:metric-decoder}--\eqref{eq:hollow-decoder}, with the deterministic
initialization and estimated-parameter profile procedure just specified,
satisfies
\begin{equation}\label{eq:dc-risk}
 \E\ell(\widehat z_D,z)\le C R_{\mathrm{prof},*}(n)
\end{equation}
for all sufficiently large $n$, with $C$ independent of $n$.
The statement is uniform over the displayed compact parameter class.
No correct grouping, true degree parameters, true block matrix, sample
splitting, or leave-one-out eigensystem is an input to the algorithm.
\end{theorem}

\begin{corollary}[The known-target-degree oracle on equal exposures]
\label{cor:known}
Under \cref{thm:dc}, suppose additionally that the actual finite-$n$ block
exposures satisfy
\begin{equation}\label{eq:equal-exposure}
 \sum_b(n_b/n)C_{ab}=S\quad\text{for all }a.
\end{equation}
Then
\[
 R_{\mathrm{prof},*}\asymp R_{\theta,*},\qquad
 \E\ell(\widehat z_D,z)\le C R_{\theta,*}.
\]
Thus there are nontrivial fixed-$K$ SBM and DC-SBM classes on which the
spectral-only fitted likelihood procedure attains a constant multiple of
the original known-parameter local LRT risk. The ordinary free-profile
decoder, with a sufficiently small floor coefficient for its structural
parameter class, also satisfies \eqref{eq:sbm-risk} for general positive
full-rank $C$ under \eqref{eq:dc-range} and \eqref{eq:equal-exposure},
without the entry-ratio-two restriction of \cref{thm:sbm}.
\end{corollary}

Exact equality in \eqref{eq:equal-exposure} can be relaxed to exposure
differences whose contribution to all likelihood contrasts is uniformly
$O(1)$. For example fixed $C\pi=S\one$ with
$n_a=n\pi_a+O(1)$ suffices. Equality of limiting exposures with no rate
is not by itself sufficient. The general known-$\theta_i$ claim without
an additional restriction is false: \cref{thm:imposs} gives an explicit
diverging risk ratio for every full-graph algorithm.

\begin{remark}[What is and is not constant]
The factor may depend on signal conditioning, group proportions, the
allowed degree-parameter range and $C_1$. It is independent of $n$ on
the stated class. We do not assert a small numerical constant, a limit of
the risk ratio, or uniformity when $\sigma_{\min}(C)\to0$. The latter
weak-signal extension needs separate quantitative estimates.
\end{remark}

